\documentclass[11pt,a4paper]{article}

\usepackage[margin=2.5cm]{geometry}
\usepackage{booktabs}
\usepackage{graphicx}
\usepackage{hyperref}
\usepackage{amsmath}
\usepackage{array}
\usepackage{xcolor}
\usepackage{enumitem}
\usepackage{caption}
\hypersetup{colorlinks=true, linkcolor=blue!60!black, urlcolor=blue!60!black}

\title{%
  \textbf{Note de synth\`ese\,: performances Charades dans VideoGLUE}\\[0.4em]
  \large Yuan et al., TMLR 2024 --- Tables~3 \`a~6\\[0.25em]
  \normalsize\url{https://arxiv.org/abs/2307.03166}
}
\author{Asma\\
  \textit{Superviseur\,: Prof.\ \'Eric Hervet}\\
  \textit{Alliance Canada --- projet Mitacs}}
\date{15~juillet~2026}

\begin{document}
\maketitle

\begin{abstract}
\noindent
Cette note r\'epond \`a la demande d'un r\'esum\'e des performances Charades
des fondation models (FMs) \'evalu\'es dans VideoGLUE (CLIP, FLAVA, CoCa,
DINOv2, VATT, VideoMAE, InternVideo), avec leurs caract\'eristiques
principales et les ressources / co\^{u}ts rapport\'es. \textbf{Oui\,:}
dans la Table~5 (MLAP), \textbf{58.2\,mAP} est le meilleur r\'esultat
Charades parmi les FMs, obtenu par \textbf{VATT}.
\end{abstract}

%==============================================================================
\section{Contexte}
%==============================================================================

VideoGLUE \cite{yuan2024videoglue} \'evalue sept FMs sur huit jeux vid\'eo
(dont Charades, classification multi-label) sous quatre m\'ethodes
d'adaptation. La Figure~1 du papier est un radar \emph{normalis\'e} (sans
chiffres bruts)\,; les scores Charades (mAP) sont dans les Tables~3--6.

\textbf{Charades (VC-ML)\,:} $\sim$30\,s / vid\'eo, plusieurs actions
simultan\'ees, $\sim$7\,811 train / 1\,814 test, 157 classes
\cite{sigurdsson2016charades}.

%==============================================================================
\section{Caract\'eristiques principales des mod\`eles}
%==============================================================================

Tous les FMs sont volontairement \'evalu\'es en \textbf{ViT-B}
\cite{dosovitskiy2021vit} pour aligner l'architecture et le co\^{u}t relatif
(Section~3.1).

\begin{table}[h]
\centering
\caption{Fondation models \'etudi\'es (d'apr\`es Table~1, VideoGLUE).}
\label{tab:fm-char}
\small
\begin{tabular}{@{}llll@{}}
\toprule
Mod\`ele & Modalit\'e pretrain & Donn\'ees & Objectif \\
\midrule
CLIP & Image + texte & WebImageText & Contrastif \\
FLAVA & Image + texte & PMD & Contrastif + MIM + MLM \\
CoCa & Image + texte & JFT-3B + ALIGN & Contrastif + captioning \\
DINOv2 & Image & LVD-142M & MIM + DINO \\
VideoMAE & Vid\'eo & Kinetics-400 & Masked video modeling \\
InternVideo & Vid\'eo & UnlabeledHybrid & MVM + contrastif \\
VATT & Vid\'eo + audio + texte & HowTo100M & Contrastif \\
\bottomrule
\end{tabular}
\end{table}

\noindent
\textbf{Image-native\,:} CLIP, FLAVA, CoCa, DINOv2.
\textbf{Video-native\,:} VideoMAE, InternVideo, VATT.

\subsection*{M\'ethodes d'adaptation (Tables~3--6)}

\begin{itemize}[nosep]
  \item \textbf{Table~3 --- End-to-end (E2E)\,:} tout le ViT-B + t\^{e}te entra\^{\i}n\'es.
  \item \textbf{Table~4 --- Frozen\,:} backbone gel\'e\,; une couche de
        cross-attention / pooler entra\^{\i}n\'ee.
  \item \textbf{Table~5 --- MLAP\,:} multi-layer attention pooler
        (plusieurs features gel\'ees, pooler hi\'erarchique entra\^{\i}n\'e).
  \item \textbf{Table~6 --- LoRA / low-rank adapter\,:} backbone gel\'e\,;
        adapters + t\^{e}te entra\^{\i}n\'es.
\end{itemize}

%==============================================================================
\section{Performances Charades (mAP)}
%==============================================================================

\begin{table}[h]
\centering
\caption{mAP Charades (\%). Chiffres extraits des Tables~3--6 de VideoGLUE.
Le meilleur FM par colonne est en gras.}
\label{tab:charades-map}
\small
\begin{tabular}{@{}lcccc@{}}
\toprule
Mod\`ele &
  Table~3 &
  Table~4 &
  Table~5 &
  Table~6 \\
 & E2E & Frozen & MLAP & LoRA \\
\midrule
CLIP        & 54.3 & 11.2 & 41.5 & 44.2 \\
FLAVA       & 48.6 & 12.6 & 38.2 & 40.8 \\
CoCa        & 55.0 &  8.8 & 19.6 & 45.8 \\
DINOv2      & 52.4 & 25.8 & 47.8 & 38.0 \\
\textbf{VATT} & \textbf{55.7} & \textbf{33.3} & \textbf{58.2} & \textbf{53.5} \\
VideoMAE    & 51.4 & 11.3 & 35.9 & 43.0 \\
InternVideo & 52.2 & 13.0 & 40.5 & 46.2 \\
\bottomrule
\end{tabular}
\end{table}

\paragraph{R\'eponse \`a la question Table~5.}
Oui\,: \textbf{58.2\,\%} est le \textbf{meilleur score Charades parmi les
sept FMs}, obtenu par \textbf{VATT} avec MLAP. C'est m\^{e}me le meilleur
score Charades FM toutes tables confondues (au-dessus de VATT E2E~55.7 et
VATT LoRA~53.5).

\paragraph{Lectures utiles.}
\begin{itemize}[nosep]
  \item En E2E, VATT et CoCa dominent Charades\,; VideoMAE atteint 51.4.
  \item Sous backbone gel\'e (Table~4), les FMs image-native s'effondrent
        souvent sur Charades\,; VATT reste nettement devant (33.3).
  \item MLAP (Table~5) r\'ecup\`ere beaucoup pour CLIP/DINOv2, et porte VATT
        \`a 58.2.
  \item Video-native FMs sont en g\'en\'eral meilleurs sur Charades quand
        le backbone est peu ou pas mis \`a jour (motion + multi-label).
\end{itemize}

%==============================================================================
\section{Ressources et co\^{u}ts}
%==============================================================================

\subsection{Ce que VideoGLUE \emph{ne} reporte \emph{pas}}

Le papier \textbf{ne publie pas} le nombre exact de GPUs, ni la m\'emoire
vid\'eo (VRAM) par carte, pour les runs des Tables~3--6. L'infra est celle
d'une \'equipe Google DeepMind\,; on peut seulement inf\'erer un entra\^{\i}nement
\textbf{multi-GPU} \`a partir du protocole (batch~256, ViT-B vid\'eo).

\subsection{Ce qui est report\'e (protocole Charades + Fig.~3)}

\begin{table}[h]
\centering
\caption{Hyperparam\`etres Charades (Appendix~D / Table~10).}
\label{tab:charades-hp}
\small
\begin{tabular}{@{}ll@{}}
\toprule
Param\`etre & Valeur \\
\midrule
Architecture & ViT-B (tous les FMs) \\
Batch size & 256 \\
Training epochs & 50 \\
Longueur de s\'equence ViT & $8 \times 14 \times 14$ tokens \\
\bottomrule
\end{tabular}
\end{table}

\begin{table}[h]
\centering
\caption{Co\^{u}ts d'adaptation approximatifs (Fig.~3)\,: param\`etres entra\^{\i}nables
et FLOPs d'inf\'erence. Les auteurs notent que, ViT-B + m\^{e}me longueur de
s\'equence $\Rightarrow$ co\^{u}ts quasi identiques entre FMs pour une m\'ethode
donn\'ee.}
\label{tab:costs}
\small
\begin{tabular}{@{}lcc@{}}
\toprule
M\'ethode d'adaptation &
  Params entra\^{\i}nables (approx.) &
  Inference FLOPs (approx.) \\
\midrule
Frozen (head seule) & $\sim$0.13\,M & $\sim$367\,B \\
MLAP & $\sim$3.7\,M & $\sim$367\,B \\
Low-rank adapter (LoRA) & $\sim$9.7\,M & $\sim$370--378\,B \\
End-to-end finetuning & $\sim$88.9\,M & $\sim$367\,B \\
\bottomrule
\end{tabular}
\end{table}

\noindent
\textbf{Nombre de param\`etres du backbone\,:} ViT-B/16 standard
$\approx$~\textbf{86\,M} param\`etres (ordre de grandeur\,; partag\'e par
les sept FMs). En E2E, $\sim$89\,M entra\^{\i}nables incluent backbone + t\^{e}te.

\noindent
\textbf{VRAM / GPUs\,:} non sp\'ecifi\'es. Un batch de~256 avec ViT-B et
clips spatio-temporels n\'ecessite typiquement une agr\'egation multi-GPU
(ou accumulation de gradients).

%==============================================================================
\begin{thebibliography}{9}

\bibitem{yuan2024videoglue}
Liangzhe Yuan et~al.
\newblock VideoGLUE: Video General Understanding Evaluation of Foundation Models.
\newblock \emph{Transactions on Machine Learning Research}, 2024.
\newblock \url{https://arxiv.org/abs/2307.03166}.

\bibitem{sigurdsson2016charades}
Gunnar~A. Sigurdsson et~al.
\newblock Hollywood in Homes: Crowdsourcing Data Collection for Activity Understanding.
\newblock \emph{ECCV}, 2016.

\bibitem{dosovitskiy2021vit}
Alexey Dosovitskiy et~al.
\newblock An Image is Worth 16x16 Words: Transformers for Image Recognition at Scale.
\newblock \emph{ICLR}, 2021.

\end{thebibliography}

\end{document}
